\documentclass[psamsfonts]{amsart}

%-------Packages---------
\usepackage{amssymb,amsfonts}
\usepackage[all,arc]{xy}
\usepackage{enumerate}
\usepackage{mathrsfs}

%--------Theorem Environments--------
%theoremstyle{plain} --- default
\newtheorem{thm}{Theorem}[section]
\newtheorem{cor}[thm]{Corollary}
\newtheorem{prop}[thm]{Proposition}
\newtheorem{lem}[thm]{Lemma}
\newtheorem{conj}[thm]{Conjecture}
\newtheorem{quest}[thm]{Question}

\theoremstyle{definition}
\newtheorem{defn}[thm]{Definition}
\newtheorem{defns}[thm]{Definitions}
\newtheorem{con}[thm]{Construction}
\newtheorem{exmp}[thm]{Example}
\newtheorem{exmps}[thm]{Examples}
\newtheorem{notn}[thm]{Notation}
\newtheorem{notns}[thm]{Notations}
\newtheorem{addm}[thm]{Addendum}
\newtheorem{exer}[thm]{Exercise}

\theoremstyle{remark}
\newtheorem{rem}[thm]{Remark}
\newtheorem{rems}[thm]{Remarks}
\newtheorem{warn}[thm]{Warning}
\newtheorem{sch}[thm]{Scholium}

\makeatletter
\let\c@equation\c@thm
\makeatother
\numberwithin{equation}{section}

\bibliographystyle{plain}

%--------Meta Data: Fill in your info------
\title{Toward $\mathcal{O}(n\log{}n)$ Large Integer Multiplication}

\author{JT Miller}

\date{DEADLINES: Draft AUGUST 19 and Final version AUGUST 31, 2019}

\begin{document}

\begin{abstract}

Large integer multiplication is useful in a number of applied fields such as cryptography as well as more theoretical fields such as numerical analysis. This paper presents an overview of large integer multiplication algorithms starting with the traditional brute force method in $\mathcal{O}(n^2)$ and moving to more sophisticated algorithms that approach the believed limit of $\mathcal{O}(n\log{}n)$. We will detail that algorithms, provide examples, and computer code for all the discussed algorithms.

\end{abstract}

\maketitle

\tableofcontents

\section{Introduction}  

Before getting to latex comments, I'll say a few words about writing, venting from many years of hard experience. 

It is important that what you write is something you would actually like to read.  There should be no overuse of symbols.
Clusters of symbols can be barbaric.  There are subjects whose literature is festooned with sentences written almost entirely with symbols like $\exists$, $\forall$.  Don't do that. Never use a symbol for a verb.  Never start a sentence with a symbol.  

Grammar should be correct.  Never start a sentence with ``Where ...".   
Mismatches of singular and plural are excruciatingly painful, utterly abhorrent. . You cannot write
``Let ..., then ..."  That is what is called a run-on sentence. You must write ``Let ... .  Then ... ."
If appropriate, ``If ..., then ..."  works just fine.

Avoid words or phrases like Simply, Obviously, Just, ..., It is easy to see.  They serve only to intimidate or to browbeat the reader into acquiescence. 
Especially if English is not your native language, make sure you have  somebody fluent in English check what you have written.     

References to numbered statements are best in the form Theorem 2.3, not just 2.3.
Equations should be referred to as (2.3), NOT 2.3 or equation 2.3 or equation (2.3).


Bad writing makes for unpleasant papers, no matter how good the material.

\subsection{History}


\section{Large Integer Multiplication Algorithms}

The table of contents command will automatically make a contents.
You must run tex at least twice for this to work.

\subsection{Traditional/Brute Force Approach}

``Environments'' are commands that are given using the \verb|\begin{}|
  and \verb|\end{}| syntax. In the preamble, you can see we've defined
  a bunch of theorem-type environments.  For example, to get a definition, 
you type:

\begin{defn}  This is how to define a definition.
\end{defn}

And for a theorem and its proof you type:

\begin{thm}
This is the statement of a theorem.
\end{thm}
\begin{proof}
And this shows that the statement is correct.
\end{proof}

Note that the numbering is taken care of automatically, and that we've predefined a bunch of these sorts of environments in the header.  These take give environments for  lemmas, corollaries and such like. 

Another useful kind of enviroment is the equation environment.  Equations
get numbered in sequence with statements, as for example

\begin{equation}  e = mc^2
\end{equation}

Note that if you do not want a numbered equation, you can use the
environment ``equation*''
 like so:

\begin{equation*}
e=mc^2
\end{equation*}
But a quicker way to tex the same command and get the same displayed result is:
\[  e=mc^2 \]

There are plenty of other equation-type enviroments that allow you to
align several equations and such like things. The AMS's guide \cite{amsshort} is a
good place to start with these.

You can also typeset math directly in a paragraph by placing it within
dollar signs.  This is called ``math mode.''  For example: Let $e$
be energy, $m$ be momentum and $c$ be the speed of light.  Then
Einstein's famous equation says that $e=mc^2$.  This is useful, but
remember that it is harder to read inline math than displayed math. 

Remember that letters get put in a different font in math mode, so
whenever you are referencing a mathematical object you should always
put it in dollar signs.  For example, $f$ is a function, but f is just
a random letter.

Both \cite{notsoshort} and \cite{amsshort} have good lists of other symbols you can use in math mode.
These include greek letters ($\alpha, \beta, \Gamma, \Delta$),
operators ($\otimes, +, \sum$) and much more ($\leq, \diamond$).

\subsection{Karatsuba and Tom Cook}

\subsection{FFT}

\subsubsection{Numerical Error}

\subsection{DWT}

\subsection{Sch\"onage-Strassen}

\subsection{Nussbaumer}

\section{Time Complexity Analysis}

If you want to draw diagrams, you should use xypic.  It's actually
much easier than it looks, and we've already included it in the header
above.  Here is an example. 
\[\xymatrix{
FX \ar[r]^-{Ff} \ar[d]_{\eta_X} & FY \ar[d]^{\eta_Y} \\
GX \ar[r]_-{Gf} & GY\\} \]

\section{Code Examples}

\section*{Acknowledgments}  
Dr. Xiang Huang was instrumental in the suggestion of this topic as a capstone project and I am deeply grateful for his tutelage, infectious energy for this and many other subjects, and the time he dedicated providing guidance and mentor-ship to an undergraduate student where there is little benefit for him. Thank you Dr. Huang. 

\begin{thebibliography}{9}

\bibitem{ams} http://www.ams.org/publications/authors/tex/amslatex

\bibitem{prime}
Crandall, Richard and Pomerance, Carl.
Prime Numbers: A Computational Perspective, 2E.
Spring Science and Business Media, Inc. 2005.

\bibitem{ada}
Goodrich, Michael T. and Tamassia, Roberto.
Algorithm Design and Applications.
John Wiley and Sons, Inc. 2015. 

\bibitem{friendly}
Silverman, Joseph H.
A Friendly Introduction to Number Theory, 4E.
Pearson Education, Inc. 2012.

\bibitem{amsshort}
Michael Downes.
Short Math Guide for \LaTeX.
http://tex.loria.fr/general/downes-short-math-guide.pdf

\bibitem{May}
J. P. May.
A Concise Course in Algebraic Topology.
University of Chicago Press. 1999. 

\bibitem{notsoshort}
Tobias Oekiter, Hubert Partl, Irene Hyna and Elisabeth Schlegl.
The Not So Short Introduction to \LaTeX 2e.
https://tobi.oetiker.ch/lshort/lshort.pdf

\end{thebibliography}

\end{document}

